\section{DLat (Latices Distributivos)}

A categoria \textbf{DLat} é fechada para todos os limites e colimites finitos. Isto significa que, a partir de latices existentes, podemos sempre construir novas estruturas que preservam a distributividade.

\subsection{Fundamentação Teórica}
As construções universais implementadas e as suas propriedades são:
\begin{itemize}
    \item \textbf{Produtos:} O produto cartesiano $L \times M$ tem as operações de \textit{meet} e \textit{join} definidas componente a componente.
    \item \textbf{Equalizadores:} O sub-latice contendo apenas os elementos onde dois homomorfismos coincidem. Como as operações são pontuais, a estrutura de latice é preservada no resultado.
    \item \textbf{Coprodutos:} Embora mais complexos, em casos finitos podem ser computados através da categoria dual de conjuntos parcialmente ordenados (Posets).
\end{itemize}

\subsection{Implementação Computacional}
Abaixo apresentam-se as funções para gerar o produto e o equalizador de latices.

\begin{lstlisting}[caption={Produtos e Equalizadores em DLat}]
% Produto de dois latices
function P = dlat_product(L, M)
    [X, Y] = meshgrid(1:length(L.elements), 1:length(M.elements));
    P.elements = [L.elements(X(:)), M.elements(Y(:))];
    P.meet = @(a, b) [L.meet(a(1), b(1)), M.meet(a(2), b(2))];
    P.join = @(a, b) [L.join(a(1), b(1)), M.join(a(2), b(2))];
end

% Equalizador de dois morfismos f e g
function E = dlat_equalizer(f, g, L)
    % Filtrar elementos onde f(x) == g(x)
    idx = arrayfun(@(i) isequal(f(L.elements{i}), g(L.elements{i})), 1:length(L.elements));
    E.elements = L.elements(idx);
    E.meet = L.meet;
    E.join = L.join;
end
\end{lstlisting}